% ===== 原文 PDF pp.9--12（印刷 pp.8--11）=====
% tr-01 · tr01a.tex · tag 区间 2.1--2.10（10 条）
% 内容：§2 标题 / §2.1 标题 / 超共形代数 (2.1)--(2.3) / 主算子 (2.4)--(2.5) / 幺正性界 (2.6) / 习题 2.1 / 缩短条件 (2.7)--(2.10)

\SEC{2}{4d $\cN=2$ 超对称基础}

% [原文 p.8]
本节在 SCFT 的语境下回顾 4d $\cN=2$ 超对称的基本概念。在 §2.1 中我们介绍 $\cN=2$ 超共形代数的幺正不可约表示；在 §2.2 中我们介绍带有两个复 Weyl 旋量超荷的超 Poincaré 代数，用 $\cN=1$ 的记号给出各超场的表达式，并写出带物质的规范理论最一般的 $\cN=2$ 拉格朗日量；库仑分支与希格斯分支的超对称真空在 §2.3 中定义；在 §2.4 中我们解释这些理论的一般非重整化定理及其在重整化群流下的行为；在 §2.5 中我们把上述拉格朗日量描述与 SCFT 联系起来。

\SUB{2.1}{$\cN=2$ 超共形对称性与表示}

我们先简要概述 4d $\cN=2$ 超共形代数及其表示。在 4d 平坦时空 $\R^{3,1}$ 中，共形代数 $\fso(4,2)$ 由洛伦兹变换 $M^\mu{}_\nu\in so(3,1)$、平移 $P^\mu$、伸缩 $D$ 以及特殊共形变换 $K_\mu$ 生成，其中 $\mu,\nu=0,1,2,3$ 是时空指标。注意伸缩 $D$ 生成阿贝尔子代数 $\fso(1,1)$，其本征值对应标度维数（权重）$\Delta$。特别地，平移生成元 $P^\mu$ 具有权重 $\Delta=1$，而特殊共形变换具有权重 $\Delta=-1$，这一点可由它们与 $D$ 的对易关系看出。关于共形代数的更多细节，我们请读者参阅 [95,96]。

共形对称性可以由 4d 时空的 Poincaré 超对称进一步扩充。\footnote{注意，能够自洽地把 Poincaré 超对称与共形对称性一起纳入的最大时空维数是 6d [97]。}当加入在洛伦兹代数 $\fso(3,1)$ 的 Weyl 表示中变换的 $\cN$ 个超荷时，共形代数 $\fso(4,2)$ 被提升为超共形代数 $\fsu(2,2|\cN)$。\footnote{当 $\cN=4$ 时，代数 $\fsu(2,2|4)$ 不是单代数，物理上相关的超共形代数由它的商 $\PSU(2,2|4)$ 给出；多出来的那个 $\fu(1)$ 在 [98,99] 中被称为 bonus 对称性。}4d $\cN=2$ 超共形代数包含如下极大玻色子子代数
\begin{equation}
  su(2,2|2)\supset so(4,2)\oplus su(2)_R\oplus u(1)_r\,,
  \tag{2.1}
\end{equation}
其中 $\fsu(2)_R\oplus u(1)_r$ 是 R 对称性代数，其生成元分别为 $R_{(mn)}$ 与 $r$，这里 $m,n=1,2$ 是 $\fsu(2)_R$ 的双重态指标。与上述共形生成元相比，超共形生成元的集合现在被 Poincaré 超荷 $\cQ^m_\alpha$、$\tilde\cQ^n_{\dot\alpha}$，连同超共形伙伴 $\cS^m_\alpha$、$\tilde\cS^n_{\dot\alpha}$，以及 R 对称性生成元所扩充。
% [Tutor Note: (2.1) 表明 $\fso(4,2)\oplus su(2)_R\oplus u(1)_r$ 是 $\fsu(2,2|2)$ 的极大玻色子子代数（原文用 $\supset$）：其中 $\fso(4,2)$ 即 4d 共形代数，多出的 $\fsu(2)_R\oplus u(1)_r$ 是 SCFT 才有的 R 对称性。另注意 $D$ 的权重约定：$[D,P^\mu]=P^\mu$、$[D,K_\mu]=-K_\mu$，故 $\Delta(P)=+1$、$\Delta(K)=-1$，$D$ 的本征值即标度维数。]

% [原文 p.9]
这里 $(\alpha,\dot\alpha)$ 表示洛伦兹代数 $\fso(3,1)$ 下的（反）手征旋量指标。注意这些超荷具有如下示意性的对易关系
\begin{equation}
  \{\cQ,\tilde\cQ\}\sim P_\mu\,,\qquad
  \{\cS,\tilde\cS\}\sim K_\mu\,,\qquad
  \{\cQ,\cS\}\sim D+M_{\mu\nu}\sigma^{\mu\nu}_{\alpha\beta}+R_{mn}+r\,,
  \tag{2.2}
\end{equation}
以及 $\{\tilde\cQ,\tilde\cS\}$ 的类似关系。它们的标度维数，以及在我们约定下 $\fu(1)_r$ 的荷，列于 (2.3)。
\begin{equation}
  \begin{array}{c|cccc}
    \hline
    & \cQ^{m}_{\alpha} & \tilde\cQ^{m}_{\dot\alpha} & \cS^{m}_{\alpha} & \tilde\cS^{m}_{\dot\alpha}\\
    \hline
    u(1)_r & \frac12 & -\frac12 & -\frac12 & \frac12\\[2pt]
    \Delta & \frac12 & \frac12 & -\frac12 & -\frac12\\
    \hline
  \end{array}
  \tag{2.3}
\end{equation}
% [Tutor Note: 读 (2.3) 时把两行分开看：$\fu(1)_r$ 行是 R 荷，$\Delta$ 行是标度维数。$\tilde\cQ$ 是 $\cQ$ 的共轭，故其 $\fu(1)_r$ 荷反号（$-\frac12$ 对 $+\frac12$）而 $\Delta$ 相同（都是 $+\frac12$）；$\cS,\tilde\cS$ 的 $\Delta=-\frac12$ 说明它们是「降」算子。$\fu(1)_r$ 荷的这套符号约定是全文的基准（表 1 的表题特意说明按此 (2.3) 对 $\mathrm{U}(1)_r$ 荷作了修改）。]

我们要强调：R 对称性属于超共形代数、从而是理论真正的对称性，这一点是 SCFT 的标志；相反地，在非共形超对称理论中它们一般并不构成对称性。

SCFT 的算子内容按超共形对称性来组织。在 $\fso(3,1)\sim su(2)\oplus su(2)$ 下具有时空量子数 $j,\bar j\in\Z$ 的算子，进一步由 $\cN=2$ 超共形代数额外 Cartan 生成元的本征值标记为 $[j,\bar j]^{(R;r)}_\Delta$，其中 $\Delta$ 是标度维数，$R\in\Z$ 是 $\fsu(2)_R$ 自旋的两倍，$r$ 标记 $\fu(1)_r$ 荷。

在幺正 SCFT 中，局域算子构成超共形代数的不可约幺正表示。与一般 CFT 一样，在平坦空间的半径量子化下，这些算子与 $S^3$ 上 Hilbert 空间中的态一一对应，并被组织成相同的表示。\footnote{更具体地说，在态--算子对应下，每个局域算子 $[j,\bar j]^{(R;r)}_\Delta$ 被等同于唯一的态 $|[j,\bar j]^{(R;r)}_\Delta\rangle$，而在本节中我们将经常省略态的 $|\cdots\rangle$ 记号。我们指出，态与局域算子之间的这种双射只有在 CFT 中才可能。}幺正性要求每个这样的表示中存在一个最低权态，称为\textbf{超共形主算子}，它满足
\begin{equation}
  \cS^m_\alpha[j,\bar j]^{(R;r)}_\Delta=0\,,\qquad
  \tilde\cS^m_{\dot\alpha}[j,\bar j]^{(R;r)}_\Delta=0\,,\qquad
  K^\mu[j,\bar j]^{(R;r)}_\Delta=0\,,
  \tag{2.4}
\end{equation}
其中最后一个等式定义了普通的共形主算子，它由前两个等式与 (2.2) 推出。给定不可约表示中的其余所有态，都可以通过把 $(\cQ,\tilde\cQ)$ 超荷作用到这些超共形主态上得到，
\begin{equation}
  \cQ^m_\alpha\cdots\tilde\cQ^n_{\dot\beta}[j,\bar j]^{(R;r)}_\Delta\,,
  \tag{2.5}
\end{equation}
% [Tutor Note: (2.4) 的三个条件并不独立。由 (2.2) 第二式 $\{\cS,\tilde\cS\}\sim K_\mu$ 立刻可得：若 $\cS$ 与 $\tilde\cS$ 都湮灭该态，则 $K_\mu$ 也湮灭它（$K_\mu$ 正比于 $\cS\tilde\cS+\tilde\cS\cS$）。这正是原文「最后一个等式由前两式与 (2.2) 推出」的含义，也是「超共形主算子必是共形主算子」的由来。]

% [原文 p.10]
因而被称为（超对称）第 $l$ 级\textbf{后裔}，其中 $l$ 计数上面出现的超荷数目。注意有些后裔同时也是共形主算子（例如一级 $(\cQ,\tilde\cQ)$ 后裔，因为对易关系 $[K,\cQ]\sim\tilde\cS$ 与 $[K,\tilde\cQ]\sim\cS$）。此外，由于超荷的费米性，给定超共形表示中共形主算子的数目是有限的。更确切地说，在具有八个实超荷的 4d $\cN=2$ SCFT 中，一个超共形表示（多重态）一般有 $2^8=64$ 个共形主算子。我们将把这些多重态称为\textbf{长多重态}。幺正性要求所有态的范数非负，从而对超共形表示给出非平凡的约束。特别地，它对超共形主算子 $\cV$ 的标度维数 $\Delta$ 施加一个由其极大玻色子子代数下的量子数决定的下界，其形式为
\begin{equation}
  \Delta\geqq\Delta_{\cV}=f(j,\bar j,R,r)\,.
  \tag{2.6}
\end{equation}
% [Tutor Note: 原文此处印作 $2^8=64$（已用 750 dpi 裁图逐字复核，确认底本确实是 $2^8=64$，不是 $2^6$）。但 $2^8=256$，与 64 并不相等，故这里存在算术上的不一致，疑为 $2^6=64$ 之误。译文照录原文的 $2^8=64$，读者按「64 个共形主算子」理解即可。]
这称为\textbf{幺正性界}。当界被饱和时，超共形表示变为可约的，并包含一个由零范数态构成的子表示，这些态称为\textbf{零态}。

\begin{exercisebox}{习题 2.1}
证明：在（超）共形代数的幺正表示中，零态自身构成一个最低权表示。
\end{exercisebox}

我们可以通过取商一致地去掉零态，所得的不可约超共形多重态与长多重态相比包含更少的算子。因此我们将这些多重态称为\textbf{短多重态}。超共形主算子为饱和幺正性界所必须满足的关系，称为\textbf{缩短条件}。鉴于 4d $\cN=2$ 超共形代数的复杂性，自然预期可以存在若干种缩短条件。首先，视构造零态所涉超荷的手征性，存在独立的手征与反手征缩短条件。因此完整的 $\cN=2$ 多重态一般由两侧的一对缩短条件指定（见表 1）。下面我们专注于手征缩短条件，它们进一步分为两类缩短类型，按 [5] 记为 $A,B$（反手征缩短类型记为 $\bar A,\bar B$）。每一缩短类型都可以有若干

% [原文 p.11]
缩短条件，相应的幺正性界可示意地写为
\begin{equation}
  \begin{aligned}
    \Delta_{\cV}&\ge f(j,\bar j,R,r)+\delta_A\,,\quad j,\bar j,R,r\ \text{不受限}\,,\\
    \Delta_{\cV}&= f(j,\bar j,R,r)+\delta_B\,,\quad j,\bar j,R,r\ \text{受限}\,,
  \end{aligned}
  \tag{2.7}
\end{equation}
其中函数 $f(j,\bar j,R,r)$ 对 $A,B$ 是同一个函数，且存在常数偏移 $\delta_{A,B}$ 满足 $\delta_A>\delta_B$。所谓不受限，是指 $A$ 可出现在任何允许的洛伦兹表示中，而缩短类型 $B$ 只在 $\cV$ 的量子数满足特定条件时才会出现。每个缩短条件还进一步由所涉超荷的手征性区分。注意类型 $A$ 与 $B$ 之间有一个重要的区别，我们现在来说明。
% [Tutor Note: (2.7) 两行共用一个编号，原始排版即如此（「一条编号跨多行时 \tag 放最后一行」）。两行的差别是本质的：类型 $A$ 是「$\ge$ + 量子数不受限」，类型 $B$ 是「$=$ + 量子数受限」。由 (2.10) 可见 $A$ 型可以在阈值处与长多重态复合（长多重态的 $\Delta$ 从上方逼近时碎裂为 $A$ 型短多重态加零态短表示）；$B$ 型因量子数取严格等式，与同量子数的其他幺正表示由有限间隙隔开，故不会发生这种复合。]

对于 (2.7) 中的类型 $A$，我们得到允许标度维数的一个下界。若界被饱和，则超共形表示将包含需要去掉的零态。在界之上，我们得到一般的长多重态（我们将记为 $L$）。因此这类短表示被称为\textbf{阈值处的短多重态}，并记为
\begin{equation}
  A_l[j,\bar j]^{(R;r)}_{\Delta_A}\,,\qquad l=1,2\,,
  \tag{2.8}
\end{equation}
其中第一个零态的形式为 $\cQ^l[j,\bar j]_{\Delta_A}$。同样重要的是要注意，零态本身构成一个短表示，而若其超共形主算子具有正范数，这样的表示将是幺正的。

对于 (2.7) 中的类型 $B$，对量子数的约束是严格等式。这意味着类型 $B$ 的短多重态由有限的间隙与具有相同洛伦兹和 R 对称性量子数的其他幺正表示隔开。这样的表示类似地记为
\begin{equation}
  B_1[j,\bar j]^{(R;r)}_{\Delta_B}\,,
  \tag{2.9}
\end{equation}
其第一个零态分别由关于 $\cQ$ 的一级后裔给出。与类型 $A$ 不同，这里去掉的零态仍构成一个子表示，但不能被提升为独立的幺正表示。

这一结构引出了\textbf{复合规则}（recombination rules）的概念。我们可以设想逐渐降低一般长多重态的标度维数 $\Delta_{\cV}$。最终它会从上方击中幺正性界，并碎裂为一个 $A$ 型短多重态加上包含所有零态的短表示 $N$（$N$ 可以是可约的）：
\begin{equation}
  L[j,\bar j]^{(R;r)}_\Delta
  \xrightarrow{\ \Delta\to\Delta_A^+\ }
  A_l[j,\bar j]^{(R;r)}_{\Delta_A}\oplus N[j_N,\bar j_N]^{(R_N;r_N)}_{\Delta_N=\Delta_A+l/2}\,.
  \tag{2.10}
\end{equation}
% [Tutor Note: 箭头上的 $\Delta\to\Delta_A^+$ 表示「从上方逼近」幺正性界：$\Delta$ 只能连续下降，故复合恰好发生在 $\Delta$ 触及界的瞬间。碎裂出的 $N$ 的维数由 $\Delta_N=\Delta_A+l/2$ 固定：$l$ 正是 (2.8) 中第一个零态所含的超荷数，每一个单位超荷带来 $1/2$ 的标度维数偏移（见 (2.3) 中 $\Delta$ 行）。]
